\chapter{Differenziazione dagli strumenti esistenti}
\label{chap:comparazione}

\section{Su quali assi ci si differenzia}
\label{sec:ch5:assi}

Nessuno degli strumenti considerati prova a fare ciò che fa questo sistema, quindi
su ciascun terreno, preso da solo, il risultato del confronto è scontato. La
differenza emerge solo tenendo insieme tre assi.

\begin{enumerate}[label=\textbf{A\arabic*.},leftmargin=*]
  \item La profondità biomeccanica, ovvero quanto il sistema arriva vicino alla
    cinematica articolare e con quale modello.
  \item L'interrogabilità, ovvero quanto si possono porre domande in linguaggio
    naturale da parte di una persona o di un agente.
  \item L'onestà dichiarata, ovvero quanto il sistema comunica, insieme al
    risultato, la qualità del dato su cui si basa.
\end{enumerate}

Gli strumenti esistenti ne coprono al massimo due, mentre lo spazio in cui si
colloca questo lavoro è proprio la combinazione dei tre.

\section{Prima famiglia, chat su metriche di allenamento}
\label{sec:ch5:allenamento}

La prima famiglia comprende le piattaforme che rendono conversazionali i dati di
allenamento. Kailo\cite{kailo}, che in precedenza si chiamava GetFast, affianca all'analisi
dell'allenamento una chat rivolta all'atleta, AI Endurance\cite{aiendurance} pubblica
un server MCP ufficiale, mentre per TrainingPeaks\cite{trainingpeaksmcp} esiste un
connettore MCP non ufficiale, segno che la domanda di
integrazione con gli agenti è già presente.

L'interrogabilità è pienamente coperta ed il motivo è lo stesso che tiene queste
piattaforme lontane dalla biomeccanica articolare, perché le grandezze esposte,
come passo, cadenza e carico, sono già aggregate ed interpretate. Non esiste
quindi, ad esempio, un angolo del ginocchio da sbagliare e per la stessa ragione il problema
dell'onestà sul dato articolare non si pone. Il sistema proposto prova invece a
portare l'interrogabilità dentro il dominio articolare.

\section{Seconda famiglia, biomeccanica markerless professionale}
\label{sec:ch5:markerless}

La seconda famiglia ricostruisce la cinematica articolare da video con
un'accuratezza validata. OpenCap\cite{opencap}, sviluppato a Stanford ed open source, lo fa a
partire da due telecamere di smartphone, mentre Theia3D\cite{theia3d} è una
soluzione commerciale a corpo intero basata su video
multicamera.

Sulla profondità biomeccanica questi sistemi sono migliori di questo lavoro 
perché stimano i punti con più precisione ed usano più telecamere e modelli
muscoloscheletrici completi. L'accuratezza di cattura dei dati di posa, però, non fa parte 
dell'ambito di questa tesi, la differenza sta negli altri due assi. Quei sistemi
producono file, report e curve costruiti per un esperto che sa leggere ciò che
riceve, senza un canale interrogabile che dichiari la propria affidabilità ad ogni
risposta. Il sistema proposto è quindi complementare, le loro uscite potrebbero
diventare una sorgente come le altre, conservate come misurate, e ricevere così
il canale interrogabile che oggi non hanno.

\section{Il protocollo non è il contributo}
\label{sec:ch5:protocollo}

Esporre degli strumenti attraverso il Model Context Protocol\cite{mcpspec} è
ormai alla portata di chiunque abbia un archivio. Un server che esponesse le
stesse venti funzioni restituendo numeri nudi sarebbe conforme quanto questo,
però non ne offrirebbe le proprietà, cioè il contratto di risposta, gli alert con
la loro fonte, il confronto di forma e la separazione dei privilegi.

\section{Tabella di comparazione}
\label{sec:ch5:sintesi}

La Tabella\ref{tab:ch5:sintesi} riassume il posizionamento di ciascuno strumento
citato. Le valutazioni riassumono gli argomenti dei paragrafi precedenti, quindi
non vanno lette come delle misure.

\begin{table}[htbp]
\centering
\caption{Posizionamento sui tre assi. Il simbolo $\bullet$ indica copertura
piena, $\circ$ copertura parziale, ``---'' assenza o non applicabilità.}
\label{tab:ch5:sintesi}
\begin{tabularx}{\textwidth}{Xlccc}
\toprule
\textbf{Strumento} & \textbf{Accesso per agenti} & \textbf{A1} & \textbf{A2} &
  \textbf{A3} \\
\midrule
Kailo (ex GetFast) & chat propria & --- & $\bullet$ & --- \\
AI Endurance & MCP ufficiale & --- & $\bullet$ & --- \\
TrainingPeaks & MCP non ufficiale & --- & $\bullet$ & --- \\
OpenCap & nessuno & $\bullet$ & --- & $\circ$ \\
Theia3D & nessuno & $\bullet$ & --- & $\circ$ \\
Questo lavoro & MCP proprietario & $\circ$ & $\bullet$ & $\bullet$ \\
\bottomrule
\end{tabularx}
\end{table}

La riga di questo lavoro porta un cerchio vuoto sulla profondità biomeccanica 
perché con i soli centri articolari alcuni piani non sono osservabili e vengono
registrati come assenze. Dichiarare una copertura piena sarebbe incoerente con
il resto della tesi.
